\documentclass[11pt, letterpaper]{article}
% ============================================================
% preamble.tex — Shared LaTeX preamble for Learning to Autolens
% ============================================================
% Contains: packages, physics notation macros, formatting.
% Include in each module via % ============================================================
% preamble.tex — Shared LaTeX preamble for Learning to Autolens
% ============================================================
% Contains: packages, physics notation macros, formatting.
% Include in each module via % ============================================================
% preamble.tex — Shared LaTeX preamble for Learning to Autolens
% ============================================================
% Contains: packages, physics notation macros, formatting.
% Include in each module via \input{../preamble}
% ============================================================

% --- Packages ---
\usepackage[utf8]{inputenc}
\usepackage[T1]{fontenc}
\usepackage{amsmath, amssymb, amsthm}
\usepackage{physics}       % \vb, \grad, \div, \curl, \dd, etc.
\usepackage{mathtools}     % \coloneqq, \prescript
\usepackage{bm}            % \bm for bold math
\usepackage{graphicx}
\usepackage{float}
\usepackage{booktabs}      % Professional tables
\usepackage{hyperref}
\usepackage{cleveref}      % \cref for smart cross-refs
\usepackage{xcolor}
\usepackage[margin=1in]{geometry}
\usepackage{enumitem}      % Custom list formatting
\usepackage{fancyhdr}      % Headers and footers
\usepackage{tcolorbox}     % Colored boxes for key results
\usepackage{listings}      % Code listings

% --- Color scheme ---
\definecolor{codeblue}{RGB}{31, 119, 180}
\definecolor{codegray}{RGB}{128, 128, 128}
\definecolor{keyresult}{RGB}{39, 123, 69}
\definecolor{exercise}{RGB}{178, 102, 0}

% --- tcolorbox environments ---
\newtcolorbox{keyresult}[1][]{
  colback=keyresult!5, colframe=keyresult!70,
  fonttitle=\bfseries, title={Key Result}, #1
}

\newtcolorbox{pythonbox}[1][]{
  colback=codeblue!5, colframe=codeblue!50,
  fonttitle=\bfseries, title={PyAutoLens Implementation}, #1
}

\newtcolorbox{exercisebox}[1][]{
  colback=exercise!5, colframe=exercise!50,
  fonttitle=\bfseries, title={Exercise}, #1
}

% --- Code listing style ---
\lstset{
  language=Python,
  basicstyle=\ttfamily\small,
  keywordstyle=\color{codeblue}\bfseries,
  commentstyle=\color{codegray}\itshape,
  stringstyle=\color{keyresult},
  breaklines=true,
  frame=single,
  framerule=0.5pt,
  rulecolor=\color{codegray!50},
}

% --- Physics macros ---
% Vectors (bold upright)
\newcommand{\vtheta}{\bm{\theta}}       % Image-plane position
\newcommand{\vbeta}{\bm{\beta}}         % Source-plane position
\newcommand{\valpha}{\bm{\alpha}}       % Deflection angle
\newcommand{\vx}{\bm{x}}               % Generic position

% Lensing quantities
\newcommand{\thetaE}{\theta_{\rm E}}    % Einstein radius
\newcommand{\SigmaCr}{\Sigma_{\rm cr}}  % Critical surface density
\newcommand{\Dd}{D_{\rm d}}             % Angular diameter distance to lens
\newcommand{\Ds}{D_{\rm s}}             % Angular diameter distance to source
\newcommand{\Dds}{D_{\rm ds}}           % Angular diameter distance lens-source

% Statistical quantities
\newcommand{\chisq}{\chi^2}             % Chi-squared
\newcommand{\chisqred}{\chi^2_{\rm red}} % Reduced chi-squared
\newcommand{\Like}{\mathcal{L}}         % Likelihood
\newcommand{\Evid}{\mathcal{Z}}         % Evidence
\newcommand{\Post}{\mathcal{P}}         % Posterior

% Abbreviations for references
\newcommand{\CK}{C\&K}                  % Congdon & Keeton (2018)
\newcommand{\NB}{N\&B}                  % Narayan & Bartelmann (1997)
\newcommand{\Night}{Nightingale+18}     % Nightingale, Dye & Massey (2018)

% --- Headers ---
\pagestyle{fancy}
\fancyhf{}
\fancyhead[L]{\small\textit{Learning to Autolens}}
\fancyhead[R]{\small\thepage}
\renewcommand{\headrulewidth}{0.4pt}

% --- Theorem environments ---
\theoremstyle{definition}
\newtheorem{definition}{Definition}[section]
\newtheorem{example}{Example}[section]

% ============================================================

% --- Packages ---
\usepackage[utf8]{inputenc}
\usepackage[T1]{fontenc}
\usepackage{amsmath, amssymb, amsthm}
\usepackage{physics}       % \vb, \grad, \div, \curl, \dd, etc.
\usepackage{mathtools}     % \coloneqq, \prescript
\usepackage{bm}            % \bm for bold math
\usepackage{graphicx}
\usepackage{float}
\usepackage{booktabs}      % Professional tables
\usepackage{hyperref}
\usepackage{cleveref}      % \cref for smart cross-refs
\usepackage{xcolor}
\usepackage[margin=1in]{geometry}
\usepackage{enumitem}      % Custom list formatting
\usepackage{fancyhdr}      % Headers and footers
\usepackage{tcolorbox}     % Colored boxes for key results
\usepackage{listings}      % Code listings

% --- Color scheme ---
\definecolor{codeblue}{RGB}{31, 119, 180}
\definecolor{codegray}{RGB}{128, 128, 128}
\definecolor{keyresult}{RGB}{39, 123, 69}
\definecolor{exercise}{RGB}{178, 102, 0}

% --- tcolorbox environments ---
\newtcolorbox{keyresult}[1][]{
  colback=keyresult!5, colframe=keyresult!70,
  fonttitle=\bfseries, title={Key Result}, #1
}

\newtcolorbox{pythonbox}[1][]{
  colback=codeblue!5, colframe=codeblue!50,
  fonttitle=\bfseries, title={PyAutoLens Implementation}, #1
}

\newtcolorbox{exercisebox}[1][]{
  colback=exercise!5, colframe=exercise!50,
  fonttitle=\bfseries, title={Exercise}, #1
}

% --- Code listing style ---
\lstset{
  language=Python,
  basicstyle=\ttfamily\small,
  keywordstyle=\color{codeblue}\bfseries,
  commentstyle=\color{codegray}\itshape,
  stringstyle=\color{keyresult},
  breaklines=true,
  frame=single,
  framerule=0.5pt,
  rulecolor=\color{codegray!50},
}

% --- Physics macros ---
% Vectors (bold upright)
\newcommand{\vtheta}{\bm{\theta}}       % Image-plane position
\newcommand{\vbeta}{\bm{\beta}}         % Source-plane position
\newcommand{\valpha}{\bm{\alpha}}       % Deflection angle
\newcommand{\vx}{\bm{x}}               % Generic position

% Lensing quantities
\newcommand{\thetaE}{\theta_{\rm E}}    % Einstein radius
\newcommand{\SigmaCr}{\Sigma_{\rm cr}}  % Critical surface density
\newcommand{\Dd}{D_{\rm d}}             % Angular diameter distance to lens
\newcommand{\Ds}{D_{\rm s}}             % Angular diameter distance to source
\newcommand{\Dds}{D_{\rm ds}}           % Angular diameter distance lens-source

% Statistical quantities
\newcommand{\chisq}{\chi^2}             % Chi-squared
\newcommand{\chisqred}{\chi^2_{\rm red}} % Reduced chi-squared
\newcommand{\Like}{\mathcal{L}}         % Likelihood
\newcommand{\Evid}{\mathcal{Z}}         % Evidence
\newcommand{\Post}{\mathcal{P}}         % Posterior

% Abbreviations for references
\newcommand{\CK}{C\&K}                  % Congdon & Keeton (2018)
\newcommand{\NB}{N\&B}                  % Narayan & Bartelmann (1997)
\newcommand{\Night}{Nightingale+18}     % Nightingale, Dye & Massey (2018)

% --- Headers ---
\pagestyle{fancy}
\fancyhf{}
\fancyhead[L]{\small\textit{Learning to Autolens}}
\fancyhead[R]{\small\thepage}
\renewcommand{\headrulewidth}{0.4pt}

% --- Theorem environments ---
\theoremstyle{definition}
\newtheorem{definition}{Definition}[section]
\newtheorem{example}{Example}[section]

% ============================================================

% --- Packages ---
\usepackage[utf8]{inputenc}
\usepackage[T1]{fontenc}
\usepackage{amsmath, amssymb, amsthm}
\usepackage{physics}       % \vb, \grad, \div, \curl, \dd, etc.
\usepackage{mathtools}     % \coloneqq, \prescript
\usepackage{bm}            % \bm for bold math
\usepackage{graphicx}
\usepackage{float}
\usepackage{booktabs}      % Professional tables
\usepackage{hyperref}
\usepackage{cleveref}      % \cref for smart cross-refs
\usepackage{xcolor}
\usepackage[margin=1in]{geometry}
\usepackage{enumitem}      % Custom list formatting
\usepackage{fancyhdr}      % Headers and footers
\usepackage{tcolorbox}     % Colored boxes for key results
\usepackage{listings}      % Code listings

% --- Color scheme ---
\definecolor{codeblue}{RGB}{31, 119, 180}
\definecolor{codegray}{RGB}{128, 128, 128}
\definecolor{keyresult}{RGB}{39, 123, 69}
\definecolor{exercise}{RGB}{178, 102, 0}

% --- tcolorbox environments ---
\newtcolorbox{keyresult}[1][]{
  colback=keyresult!5, colframe=keyresult!70,
  fonttitle=\bfseries, title={Key Result}, #1
}

\newtcolorbox{pythonbox}[1][]{
  colback=codeblue!5, colframe=codeblue!50,
  fonttitle=\bfseries, title={PyAutoLens Implementation}, #1
}

\newtcolorbox{exercisebox}[1][]{
  colback=exercise!5, colframe=exercise!50,
  fonttitle=\bfseries, title={Exercise}, #1
}

% --- Code listing style ---
\lstset{
  language=Python,
  basicstyle=\ttfamily\small,
  keywordstyle=\color{codeblue}\bfseries,
  commentstyle=\color{codegray}\itshape,
  stringstyle=\color{keyresult},
  breaklines=true,
  frame=single,
  framerule=0.5pt,
  rulecolor=\color{codegray!50},
}

% --- Physics macros ---
% Vectors (bold upright)
\newcommand{\vtheta}{\bm{\theta}}       % Image-plane position
\newcommand{\vbeta}{\bm{\beta}}         % Source-plane position
\newcommand{\valpha}{\bm{\alpha}}       % Deflection angle
\newcommand{\vx}{\bm{x}}               % Generic position

% Lensing quantities
\newcommand{\thetaE}{\theta_{\rm E}}    % Einstein radius
\newcommand{\SigmaCr}{\Sigma_{\rm cr}}  % Critical surface density
\newcommand{\Dd}{D_{\rm d}}             % Angular diameter distance to lens
\newcommand{\Ds}{D_{\rm s}}             % Angular diameter distance to source
\newcommand{\Dds}{D_{\rm ds}}           % Angular diameter distance lens-source

% Statistical quantities
\newcommand{\chisq}{\chi^2}             % Chi-squared
\newcommand{\chisqred}{\chi^2_{\rm red}} % Reduced chi-squared
\newcommand{\Like}{\mathcal{L}}         % Likelihood
\newcommand{\Evid}{\mathcal{Z}}         % Evidence
\newcommand{\Post}{\mathcal{P}}         % Posterior

% Abbreviations for references
\newcommand{\CK}{C\&K}                  % Congdon & Keeton (2018)
\newcommand{\NB}{N\&B}                  % Narayan & Bartelmann (1997)
\newcommand{\Night}{Nightingale+18}     % Nightingale, Dye & Massey (2018)

% --- Headers ---
\pagestyle{fancy}
\fancyhf{}
\fancyhead[L]{\small\textit{Learning to Autolens}}
\fancyhead[R]{\small\thepage}
\renewcommand{\headrulewidth}{0.4pt}

% --- Theorem environments ---
\theoremstyle{definition}
\newtheorem{definition}{Definition}[section]
\newtheorem{example}{Example}[section]


\fancyhead[C]{\small\textbf{Module 01: Grids, Galaxies, and Ray-Tracing}}

\title{\textbf{Module 01: Grids, Galaxies, and Ray-Tracing}\\[0.5em]
\large Theory Companion — Learning to Autolens}
\author{Rodrigo C\'ordova Rosado\\
\small Harvard-Smithsonian Center for Astrophysics\\
\small \texttt{rodrigo.cordova\_rosado@cfa.harvard.edu}}
\date{\today}

\begin{document}
\maketitle

\begin{abstract}
This document provides the mathematical background for Module~01 of
\emph{Learning to Autolens}. We derive the gravitational lens equation,
define the key lensing quantities (convergence, shear, magnification),
and describe the mass and light profile models used in PyAutoLens.
All results are self-contained; companion references are
Congdon \& Keeton (2018, hereafter \CK), Narayan \& Bartelmann (1997,
hereafter \NB), and Schneider, Ehlers \& Falco (1992, hereafter S92).
\end{abstract}

\tableofcontents
\newpage

% ============================================================
\section{The Gravitational Lens Equation}
\label{sec:lens-equation}
% ============================================================

The fundamental equation of gravitational lensing maps angular positions
in the \textbf{image plane} $\vtheta$ to the \textbf{source plane}
$\vbeta$ via the scaled deflection angle $\valpha$:

\begin{keyresult}
\begin{equation}
\vbeta = \vtheta - \valpha(\vtheta)
\label{eq:lens-equation}
\end{equation}
\end{keyresult}

\noindent where $\vtheta$ and $\vbeta$ are two-dimensional angular
positions (in arc-seconds), and $\valpha(\vtheta)$ is the scaled
deflection angle at position $\vtheta$ (\CK\ eq.~4.1; \NB\ eq.~3).

The deflection angle is related to the physical deflection
$\hat{\valpha}$ by the distance ratio:
\begin{equation}
\valpha = \frac{\Dds}{\Ds} \hat{\valpha}
\label{eq:scaled-deflection}
\end{equation}
where $\Dd$, $\Ds$, and $\Dds$ are the angular diameter distances to
the lens, to the source, and between the lens and source, respectively.

\subsection{Critical Surface Density}

The strength of lensing is governed by the \textbf{critical surface
mass density}:
\begin{equation}
\SigmaCr = \frac{c^2}{4\pi G} \frac{\Ds}{\Dd \Dds}
\label{eq:sigma-cr}
\end{equation}
(\CK\ eq.~4.5; \NB\ eq.~10). This sets the scale: a lens with
projected surface density $\Sigma \geq \SigmaCr$ produces multiple
images (strong lensing).

\subsection{Einstein Radius}

The \textbf{Einstein radius} is the characteristic angular scale of
strong lensing:
\begin{equation}
\thetaE = \sqrt{\frac{4GM(\thetaE)}{c^2} \frac{\Dds}{\Dd \Ds}}
\label{eq:einstein-radius}
\end{equation}
where $M(\thetaE)$ is the projected mass enclosed within $\thetaE$
(\CK\ eq.~4.9). For a typical galaxy-scale lens,
$\thetaE \sim 0.5''$--$3''$.

% ============================================================
\section{Convergence, Shear, and Magnification}
\label{sec:convergence-shear}
% ============================================================

\subsection{Convergence}

The \textbf{convergence} is the dimensionless projected surface mass density:
\begin{equation}
\kappa(\vtheta) = \frac{\Sigma(\vtheta)}{\SigmaCr}
\label{eq:convergence}
\end{equation}

\subsection{Lensing Potential}

The deflection angle derives from a 2D potential $\psi(\vtheta)$:
\begin{equation}
\valpha = \nabla \psi, \qquad \nabla^2 \psi = 2\kappa
\label{eq:lensing-potential}
\end{equation}
(\CK\ eqs.~4.14--4.15; \NB\ eq.~8). The 2D Poisson equation
$\nabla^2 \psi = 2\kappa$ is the central equation relating mass to
deflection.

\subsection{Jacobian Matrix and Magnification}

The lens mapping $\vtheta \to \vbeta$ has the Jacobian matrix:
\begin{equation}
\mathbf{A} = \frac{\partial \vbeta}{\partial \vtheta}
= \delta_{ij} - \frac{\partial^2 \psi}{\partial \theta_i \partial \theta_j}
= \begin{pmatrix}
1 - \kappa - \gamma_1 & -\gamma_2 \\
-\gamma_2 & 1 - \kappa + \gamma_1
\end{pmatrix}
\label{eq:jacobian}
\end{equation}
where $\gamma_1$ and $\gamma_2$ are the components of the \textbf{shear}
$\gamma = \gamma_1 + i\gamma_2$ (\CK\ eq.~5.7; \NB\ eqs.~17--19).

The \textbf{magnification} is the inverse determinant:
\begin{keyresult}
\begin{equation}
\mu = \frac{1}{\det \mathbf{A}} = \frac{1}{(1-\kappa)^2 - |\gamma|^2}
\label{eq:magnification}
\end{equation}
\end{keyresult}

\subsection{Critical Curves and Caustics}

\textbf{Critical curves} are the loci in the image plane where
$\det \mathbf{A} = 0$ (infinite magnification). There are two types:
\begin{align}
\text{Tangential:}\quad & 1 - \kappa - |\gamma| = 0 \label{eq:tangential-cc} \\
\text{Radial:}\quad     & 1 - \kappa + |\gamma| = 0 \label{eq:radial-cc}
\end{align}

The image of a critical curve under the lens mapping $\vtheta \to \vbeta$
is called a \textbf{caustic}. When a source crosses a caustic, images are
created or destroyed in pairs (\CK\ Ch.~5; Petters et al.\ 2001).

% ============================================================
\section{Light Profiles: The S\'ersic Function}
\label{sec:sersic}
% ============================================================

Galaxy surface brightness is described by the \textbf{S\'ersic profile}
(S\'ersic 1963):
\begin{equation}
I(R) = I_e \exp\left\{ -b_n \left[ \left(\frac{R}{R_e}\right)^{1/n} - 1 \right] \right\}
\label{eq:sersic}
\end{equation}
where $R_e$ is the effective (half-light) radius, $n$ is the S\'ersic
index, and $b_n \approx 2n - 1/3 + 4/(405n)$ ensures that $R_e$
encloses half the total light (Ciotti \& Bertin 1999).

Special cases:
\begin{itemize}[nosep]
  \item $n = 1$: exponential profile (disk galaxies)
  \item $n = 4$: de Vaucouleurs profile (elliptical galaxies)
\end{itemize}

\subsection{Ellipticity Parameterization}

PyAutoLens parameterizes ellipticity using components
$(\epsilon_1, \epsilon_2)$ rather than axis ratio $q$ and position
angle $\phi$:
\begin{equation}
\epsilon_1 = \frac{1-q}{1+q} \sin(2\phi), \qquad
\epsilon_2 = \frac{1-q}{1+q} \cos(2\phi)
\label{eq:ell-comps}
\end{equation}
with total ellipticity $|\epsilon| = (1-q)/(1+q)$. This avoids
discontinuities at $\phi = 0, \pi$ and is standard in weak lensing
(\NB\ Sec.~4.1).

% ============================================================
\section{Mass Profiles}
\label{sec:mass-profiles}
% ============================================================

\subsection{Singular Isothermal Sphere (SIS)}

The SIS has density $\rho(r) = \sigma_v^2 / (2\pi G r^2)$, giving:
\begin{align}
\kappa(\theta) &= \frac{\thetaE}{2\theta} \label{eq:sis-kappa} \\
\alpha(\theta) &= \thetaE \quad \text{(constant!)} \label{eq:sis-alpha}
\end{align}
with Einstein radius (\CK\ eq.~4.28):
\begin{equation}
\thetaE = 4\pi \left(\frac{\sigma_v}{c}\right)^2 \frac{\Dds}{\Ds}
\label{eq:sis-thetaE}
\end{equation}

The mean convergence inside $\thetaE$ is $\langle\kappa\rangle = 1$,
which is the defining property of the Einstein radius.

\subsection{Singular Isothermal Ellipsoid (SIE)}

The SIE generalizes the SIS to elliptical isodensity contours
(Kormann, Schneider \& Bartelmann 1994). It is the workhorse model for
galaxy-scale lenses because it produces the correct image configurations
(doubles and quads) with physically motivated mass distributions.

In PyAutoLens, the SIE is \texttt{al.mp.Isothermal} with parameters:
\texttt{centre}, \texttt{ell\_comps}, \texttt{einstein\_radius}.

\subsection{NFW Profile}

The Navarro-Frenk-White (1996) profile describes dark matter halos:
\begin{equation}
\rho(r) = \frac{\rho_s}{(r/r_s)(1 + r/r_s)^2}
\label{eq:nfw-density}
\end{equation}
with known analytic convergence (Bartelmann 1996; \CK\ Sec.~4.5).

% ============================================================
\section{Multi-Plane Lensing}
\label{sec:multiplane}
% ============================================================

For mass at multiple redshifts, light is deflected sequentially at each
lens plane. The multi-plane lens equation (S92, Ch.~9):
\begin{equation}
\vtheta_{j+1} = \vtheta_1 - \sum_{i=1}^{j} \beta_{ij} \valpha_i(\vtheta_i)
\label{eq:multiplane}
\end{equation}
where $\beta_{ij} = D_{ij} D_s / (D_j D_{is})$ are the distance ratios.
PyAutoLens handles this automatically via the \texttt{Tracer} object.

% ============================================================
\section{Cosmological Distances}
\label{sec:distances}
% ============================================================

All lensing distances are \textbf{angular diameter distances} in a
Friedmann-Lema\^itre-Robertson-Walker cosmology:
\begin{equation}
D_A(z_1, z_2) = \frac{1}{1+z_2} \int_{z_1}^{z_2} \frac{c\,dz'}{H(z')}
\label{eq:angular-diameter-distance}
\end{equation}
for a flat universe, where $H(z) = H_0\sqrt{\Omega_m(1+z)^3 + \Omega_\Lambda}$.

\begin{pythonbox}
PyAutoLens wraps \texttt{astropy} cosmology:
\begin{lstlisting}
cosmo = al.cosmo.Planck15()
D_d = cosmo.angular_diameter_distance_to_earth_in_kpc_from(redshift=0.5)
Sigma_cr = cosmo.critical_surface_density_between_redshifts_solar_mass_per_kpc2_from(
    redshift_0=0.5, redshift_1=1.0
)
\end{lstlisting}
\end{pythonbox}

% ============================================================
\section*{References}
% ============================================================

\begin{itemize}[nosep, leftmargin=*]
  \item Bartelmann, M.\ (1996), A\&A, 313, 697 — NFW lensing properties
  \item Ciotti, L.\ \& Bertin, G.\ (1999), A\&A, 352, 447 — S\'ersic $b_n$ approximation
  \item Congdon, A.B.\ \& Keeton, C.R.\ (2018), \emph{Principles of Gravitational Lensing}, Springer
  \item Kormann, R., Schneider, P.\ \& Bartelmann, M.\ (1994), A\&A, 284, 285 — SIE model
  \item Narayan, R.\ \& Bartelmann, M.\ (1997), arXiv:astro-ph/9606001
  \item Navarro, J.F., Frenk, C.S.\ \& White, S.D.M.\ (1996), ApJ, 462, 563
  \item Nightingale, J.W., Dye, S.\ \& Massey, R.J.\ (2018), MNRAS, 478, 4738
  \item Petters, A.O., Levine, H.\ \& Wambsganss, J.\ (2001), \emph{Singularity Theory and Gravitational Lensing}, Birkh\"auser
  \item Schneider, P., Ehlers, J.\ \& Falco, E.E.\ (1992), \emph{Gravitational Lenses}, Springer
  \item S\'ersic, J.L.\ (1963), Bolet\'in de la Asociaci\'on Argentina de Astronom\'ia, 6, 41
\end{itemize}

\end{document}
